% Options for packages loaded elsewhere
\PassOptionsToPackage{unicode}{hyperref}
\PassOptionsToPackage{hyphens}{url}
\documentclass[
]{article}
\usepackage{xcolor}
\usepackage{amsmath,amssymb}
\setcounter{secnumdepth}{-\maxdimen} % remove section numbering
\usepackage{iftex}
\ifPDFTeX
  \usepackage[T1]{fontenc}
  \usepackage[utf8]{inputenc}
  \usepackage{textcomp} % provide euro and other symbols
\else % if luatex or xetex
  \usepackage{unicode-math} % this also loads fontspec
  \defaultfontfeatures{Scale=MatchLowercase}
  \defaultfontfeatures[\rmfamily]{Ligatures=TeX,Scale=1}
\fi
\usepackage{lmodern}
\ifPDFTeX\else
  % xetex/luatex font selection
\fi
% Use upquote if available, for straight quotes in verbatim environments
\IfFileExists{upquote.sty}{\usepackage{upquote}}{}
\IfFileExists{microtype.sty}{% use microtype if available
  \usepackage[]{microtype}
  \UseMicrotypeSet[protrusion]{basicmath} % disable protrusion for tt fonts
}{}
\makeatletter
\@ifundefined{KOMAClassName}{% if non-KOMA class
  \IfFileExists{parskip.sty}{%
    \usepackage{parskip}
  }{% else
    \setlength{\parindent}{0pt}
    \setlength{\parskip}{6pt plus 2pt minus 1pt}}
}{% if KOMA class
  \KOMAoptions{parskip=half}}
\makeatother
\usepackage{color}
\usepackage{fancyvrb}
\newcommand{\VerbBar}{|}
\newcommand{\VERB}{\Verb[commandchars=\\\{\}]}
\DefineVerbatimEnvironment{Highlighting}{Verbatim}{commandchars=\\\{\}}
% Add ',fontsize=\small' for more characters per line
\newenvironment{Shaded}{}{}
\newcommand{\AlertTok}[1]{\textcolor[rgb]{1.00,0.00,0.00}{\textbf{#1}}}
\newcommand{\AnnotationTok}[1]{\textcolor[rgb]{0.38,0.63,0.69}{\textbf{\textit{#1}}}}
\newcommand{\AttributeTok}[1]{\textcolor[rgb]{0.49,0.56,0.16}{#1}}
\newcommand{\BaseNTok}[1]{\textcolor[rgb]{0.25,0.63,0.44}{#1}}
\newcommand{\BuiltInTok}[1]{\textcolor[rgb]{0.00,0.50,0.00}{#1}}
\newcommand{\CharTok}[1]{\textcolor[rgb]{0.25,0.44,0.63}{#1}}
\newcommand{\CommentTok}[1]{\textcolor[rgb]{0.38,0.63,0.69}{\textit{#1}}}
\newcommand{\CommentVarTok}[1]{\textcolor[rgb]{0.38,0.63,0.69}{\textbf{\textit{#1}}}}
\newcommand{\ConstantTok}[1]{\textcolor[rgb]{0.53,0.00,0.00}{#1}}
\newcommand{\ControlFlowTok}[1]{\textcolor[rgb]{0.00,0.44,0.13}{\textbf{#1}}}
\newcommand{\DataTypeTok}[1]{\textcolor[rgb]{0.56,0.13,0.00}{#1}}
\newcommand{\DecValTok}[1]{\textcolor[rgb]{0.25,0.63,0.44}{#1}}
\newcommand{\DocumentationTok}[1]{\textcolor[rgb]{0.73,0.13,0.13}{\textit{#1}}}
\newcommand{\ErrorTok}[1]{\textcolor[rgb]{1.00,0.00,0.00}{\textbf{#1}}}
\newcommand{\ExtensionTok}[1]{#1}
\newcommand{\FloatTok}[1]{\textcolor[rgb]{0.25,0.63,0.44}{#1}}
\newcommand{\FunctionTok}[1]{\textcolor[rgb]{0.02,0.16,0.49}{#1}}
\newcommand{\ImportTok}[1]{\textcolor[rgb]{0.00,0.50,0.00}{\textbf{#1}}}
\newcommand{\InformationTok}[1]{\textcolor[rgb]{0.38,0.63,0.69}{\textbf{\textit{#1}}}}
\newcommand{\KeywordTok}[1]{\textcolor[rgb]{0.00,0.44,0.13}{\textbf{#1}}}
\newcommand{\NormalTok}[1]{#1}
\newcommand{\OperatorTok}[1]{\textcolor[rgb]{0.40,0.40,0.40}{#1}}
\newcommand{\OtherTok}[1]{\textcolor[rgb]{0.00,0.44,0.13}{#1}}
\newcommand{\PreprocessorTok}[1]{\textcolor[rgb]{0.74,0.48,0.00}{#1}}
\newcommand{\RegionMarkerTok}[1]{#1}
\newcommand{\SpecialCharTok}[1]{\textcolor[rgb]{0.25,0.44,0.63}{#1}}
\newcommand{\SpecialStringTok}[1]{\textcolor[rgb]{0.73,0.40,0.53}{#1}}
\newcommand{\StringTok}[1]{\textcolor[rgb]{0.25,0.44,0.63}{#1}}
\newcommand{\VariableTok}[1]{\textcolor[rgb]{0.10,0.09,0.49}{#1}}
\newcommand{\VerbatimStringTok}[1]{\textcolor[rgb]{0.25,0.44,0.63}{#1}}
\newcommand{\WarningTok}[1]{\textcolor[rgb]{0.38,0.63,0.69}{\textbf{\textit{#1}}}}
\usepackage{longtable,booktabs,array}
\newcounter{none} % for unnumbered tables
\usepackage{calc} % for calculating minipage widths
% Correct order of tables after \paragraph or \subparagraph
\usepackage{etoolbox}
\makeatletter
\patchcmd\longtable{\par}{\if@noskipsec\mbox{}\fi\par}{}{}
\makeatother
% Allow footnotes in longtable head/foot
\IfFileExists{footnotehyper.sty}{\usepackage{footnotehyper}}{\usepackage{footnote}}
\makesavenoteenv{longtable}
\setlength{\emergencystretch}{3em} % prevent overfull lines
\providecommand{\tightlist}{%
  \setlength{\itemsep}{0pt}\setlength{\parskip}{0pt}}
\usepackage{bookmark}
\IfFileExists{xurl.sty}{\usepackage{xurl}}{} % add URL line breaks if available
\urlstyle{same}
\hypersetup{
  hidelinks,
  pdfcreator={LaTeX via pandoc}}

\author{}
\date{}

\begin{document}

\section{Multi-Modal Stock Price Prediction Using LSTM with Technical
Indicators and News
Sentiment}\label{multi-modal-stock-price-prediction-using-lstm-with-technical-indicators-and-news-sentiment}

\textbf{Authors:} Adarsh Singh, Venkatesh Nagarjuna, Mayur Patil\\
\textbf{Course:} DATS 6303 - Deep Learning (Fall 2025)\\
\textbf{Instructor:} Dr. Amir Jafari\\
\textbf{GitHub Repository:}
\url{https://github.com/drsh0755/FinalProject-Group5}

\begin{center}\rule{0.5\linewidth}{0.5pt}\end{center}

\subsection{Abstract}\label{abstract}

Stock price prediction remains a challenging problem in financial
markets due to the complex interplay between quantitative patterns and
qualitative information. This project presents a deep learning approach
to predict SPY (S\&P 500 ETF) and QQQ closing prices by integrating
technical indicators with financial news sentiment using a Long
Short-Term Memory (LSTM) neural network implemented in PyTorch. The
model combines 39 technical indicators derived from price and volume
data with 5 sentiment scores extracted from financial news articles. We
employ comprehensive regularization techniques including dropout, batch
normalization, and early stopping to prevent overfitting. The model
architecture includes 2-layer stacked LSTM with 64 hidden units, batch
normalization layers, and gradient clipping for stable training. The
results demonstrate that incorporating news sentiment alongside
technical analysis provides meaningful predictive capability for stock
price forecasting.

\begin{center}\rule{0.5\linewidth}{0.5pt}\end{center}

\subsection{1. Introduction}\label{1-introduction}

\subsubsection{1.1 Motivation and Problem
Statement}\label{11-motivation-and-problem-statement}

Financial markets are influenced by a complex combination of historical
price patterns and real-time information from news and events.
Traditional technical analysis focuses exclusively on price history,
while sentiment analysis in isolation may miss critical momentum
patterns. This project addresses the stock price prediction problem by
fusing technical indicators with financial news sentiment in a unified
deep learning framework.

The primary objective is to predict the closing price of SPY (SPDR S\&P
500 ETF Trust) using a PyTorch-based LSTM model that processes sequences
of multi-modal features. The SPY ETF tracks the S\&P 500 index and
serves as a broad market benchmark, making it an ideal candidate for
validating predictive models due to its high liquidity and extensive
news coverage.

\subsubsection{1.2 Approach Overview}\label{12-approach-overview}

Our approach employs a two-stage data pipeline followed by LSTM-based
modeling:

\begin{enumerate}
\def\labelenumi{\arabic{enumi}.}
\item
  \textbf{Feature Engineering Stage:} We construct a rich feature set
  combining 39 technical indicators with daily sentiment scores
  extracted from financial news articles.
\item
  \textbf{Deep Learning Stage:} A 2-layer stacked LSTM network with
  batch normalization and dropout processes 30-day sequences of these
  features to predict the next day\textquotesingle s closing price.
\end{enumerate}

The model is trained using Mean Squared Error (MSE) loss with Adam
optimization, learning rate scheduling, gradient clipping, and early
stopping to ensure generalization. Unlike simpler approaches that treat
stock prediction as a univariate time series problem, our multi-modal
architecture captures both market dynamics and information shocks from
news events.

\begin{center}\rule{0.5\linewidth}{0.5pt}\end{center}

\subsection{2. Data preparation
pipeline}\label{2-data-preparation-pipeline}

\subsubsection{2.1 Data Sources}\label{21-data-sources}

The dataset integrates two primary data streams:

\textbf{Technical Analysis Data (Stream 1):} Historical stock price and
volume data for SPY was obtained from Yahoo Finance using the
\texttt{yfinance} Python library. The raw data includes daily Open,
High, Low, Close (OHLC) prices and trading Volume.

\textbf{News Sentiment Data (Stream 2):} Financial news articles related
to SPY were collected from the Alpha Vantage News Sentiment API. Each
article includes a pre-computed sentiment score ranging from -1
(negative) to +1 (positive), along with a relevance score indicating the
article\textquotesingle s pertinence to the specified ticker.

\subsubsection{2.2 Dataset
Characteristics}\label{22-dataset-characteristics}

{\def\LTcaptype{none} % do not increment counter
\begin{longtable}[]{@{}lll@{}}
\toprule\noalign{}
Component & Source & Details \\
\midrule\noalign{}
\endhead
\bottomrule\noalign{}
\endlastfoot
\textbf{Stock Price Data} & Yahoo Finance (yfinance) & SPY (SPDR S\&P
500 ETF Trust) historical daily OHLCV data \\
\textbf{News Sentiment Data} & Alpha Vantage News API & 6,700 financial
news articles with sentiment scores (October 2024 -- December 2025) \\
\textbf{Date Range} & April 16, 2024 -- November 26, 2025 & 407 trading
days of historical data \\
\textbf{Total Records} & 407 trading days & After preprocessing and
alignment \\
\end{longtable}
}

\subsubsection{2.3 Dataset Statistics}\label{23-dataset-statistics}

{\def\LTcaptype{none} % do not increment counter
\begin{longtable}[]{@{}ll@{}}
\toprule\noalign{}
Metric & Value \\
\midrule\noalign{}
\endhead
\bottomrule\noalign{}
\endlastfoot
\textbf{Total Trading Days} & 407 \\
\textbf{Total News Articles} & 6,700 \\
\textbf{30-Day Sliding Sequences} & 378 \\
\textbf{Training Sequences (70\%)} & 264 \\
\textbf{Validation Sequences (15\%)} & 56 \\
\textbf{Test Sequences (15\%)} & 58 \\
\textbf{Total Input Features} & 44 \\
\textbf{Technical Indicators} & 39 features \\
\textbf{Sentiment Features} & 5 features \\
\end{longtable}
}

\subsubsection{2.4 News Sentiment
Statistics}\label{24-news-sentiment-statistics}

Financial sentiment aggregation across the 6,700 articles provides
nuanced market sentiment:

{\def\LTcaptype{none} % do not increment counter
\begin{longtable}[]{@{}ll@{}}
\toprule\noalign{}
Metric & Value \\
\midrule\noalign{}
\endhead
\bottomrule\noalign{}
\endlastfoot
\textbf{Total Articles Indexed} & 6,700 \\
\textbf{Sentiment Score Range} & 0.001 to 0.352 \\
\textbf{Mean Sentiment} & 0.056 \\
\textbf{Sentiment Features} & Mean, Median, Std Dev, Min, Max (per
trading day) \\
\end{longtable}
}

\begin{center}\rule{0.5\linewidth}{0.5pt}\end{center}

\subsubsection{2.5 Feature Engineering}\label{25-feature-engineering}

\paragraph{\texorpdfstring{\textbf{Technical Indicators (39
Features)}}{Technical Indicators (39 Features)}}\label{technical-indicators-39-features}

The system extracts comprehensive technical indicators across multiple
categories:

\textbf{Moving Averages (8 features)}

\begin{itemize}
\item
  Simple Moving Averages (SMA): 5-day, 10-day, 20-day, 50-day
\item
  Exponential Moving Averages (EMA): 5-day, 10-day, 20-day, 50-day
\end{itemize}

\textbf{Momentum Indicators (7 features)}

\begin{itemize}
\item
  Relative Strength Index (RSI)
\item
  MACD, MACD Signal, MACD Histogram
\item
  Stochastic Oscillator
\item
  Williams \%R
\end{itemize}

\textbf{Volatility Indicators (6 features)}

\begin{itemize}
\item
  Bollinger Bands: Upper, Middle, Lower, Width, Position
\item
  Average True Range (ATR)
\item
  10-day, 20-day, 50-day Volatility
\end{itemize}

\textbf{Volume Indicators (4 features)}

\begin{itemize}
\item
  On-Balance Volume (OBV)
\item
  Volume SMA (20-day)
\item
  Volume Ratio (current volume / SMA)
\item
  Daily Volume Rate of Change
\end{itemize}

\textbf{Trend Indicators (2 features)}

\begin{itemize}
\item
  Average Directional Index (ADX)
\item
  Commodity Channel Index (CCI)
\end{itemize}

\textbf{Lagged Features (8 features)}

\begin{itemize}
\item
  Close price lags: 1-day, 2-day, 3-day, 5-day
\item
  Return lags: 1-day, 2-day, 3-day, 5-day
\end{itemize}

\textbf{Price Action Features (4 features)}

\begin{itemize}
\item
  Daily returns, Log returns
\item
  High-Low percentage, Open-Close percentage
\item
  Price position (normalized position within daily range)
\end{itemize}

\paragraph{Sentiment Features (5
Features)}\label{sentiment-features-5-features}

Daily sentiment aggregation from news articles:

\begin{itemize}
\item
  \textbf{Sentiment Mean}: Average sentiment score across articles for
  the trading day
\item
  \textbf{Sentiment Median}: Median sentiment score
\item
  \textbf{Sentiment Std Dev}: Volatility of sentiment opinions
\item
  \textbf{Sentiment Min/Max}: Extreme sentiment bounds
\end{itemize}

\textbf{Aggregation Method}: All news articles for each trading day are
aggregated using statistical measures, capturing both central tendency
and dispersion of market sentiment.

\subsubsection{2.6 Data Preprocessing}\label{26-data-preprocessing}

The preprocessing pipeline consists of the following steps:

\begin{enumerate}
\def\labelenumi{\arabic{enumi}.}
\item
  \textbf{Date Alignment:} Stock price data and news sentiment data were
  merged on trading dates, ensuring temporal consistency.
\item
  \textbf{Missing Value Handling:} NaN values resulting from indicator
  computation windows were removed. Infinite values were replaced with
  zeros.
\item
  \textbf{Normalization:} All features and the target variable (closing
  price) were scaled to the range using MinMaxScaler from scikit-learn.
  Separate scalers were fitted on the training set and applied to
  validation and test sets to prevent data leakage.
\item
  \textbf{Sequence Creation:} Time series sequences of length 30 were
  constructed using a sliding window approach. Each sequence contains 30
  consecutive days of 43 features, with the target being the closing
  price on day 31.
\end{enumerate}

\begin{center}\rule{0.5\linewidth}{0.5pt}\end{center}

\subsection{3. Description of the Deep Learning Network and Training
Algorithm}\label{3-description-of-the-deep-learning-network-and-training-algorithm}

\subsubsection{3.1 Network Architecture}\label{31-network-architecture}

The model is implemented as a custom PyTorch neural network class that
extends \texttt{nn.Module}. The architecture consists of three major
components:

\paragraph{3.1.1 LSTM Layers}\label{311-lstm-layers}

The core of the model is a 2-layer stacked LSTM that processes
sequential input data:

\begin{Shaded}
\begin{Highlighting}[]
\VariableTok{self}\NormalTok{.lstm }\OperatorTok{=}\NormalTok{ nn.LSTM(}
\NormalTok{    input\_size}\OperatorTok{=}\DecValTok{43}\NormalTok{,}
\NormalTok{    hidden\_size}\OperatorTok{=}\DecValTok{64}\NormalTok{,}
\NormalTok{    num\_layers}\OperatorTok{=}\DecValTok{2}\NormalTok{,}
\NormalTok{    dropout}\OperatorTok{=}\FloatTok{0.4}\NormalTok{,}
\NormalTok{    batch\_first}\OperatorTok{=}\VariableTok{True}
\NormalTok{)}
\end{Highlighting}
\end{Shaded}

\textbf{Configuration:}

\begin{itemize}
\item
  Input size: 43 features per time step
\item
  Hidden size: 64 units per LSTM layer
\item
  Number of layers: 2 stacked LSTM layers
\item
  Inter-layer dropout: 0.4
\item
  Batch-first format: Input shape (batch\_size, sequence\_length,
  input\_size)
\end{itemize}

The LSTM layers capture temporal dependencies and patterns in the 30-day
sequences. The hidden state dimensionality of 64 was chosen to balance
model capacity and generalization, reduced from larger values to prevent
overfitting on the limited dataset.

\paragraph{3.1.2 Fully Connected Layers with
Regularization}\label{312-fully-connected-layers-with-regularization}

After the LSTM processes the sequence, the final hidden state is passed
through a series of fully connected layers with batch normalization and
dropout:

\begin{Shaded}
\begin{Highlighting}[]
\VariableTok{self}\NormalTok{.fc1 }\OperatorTok{=}\NormalTok{ nn.Linear(}\DecValTok{64}\NormalTok{, }\DecValTok{32}\NormalTok{)}
\VariableTok{self}\NormalTok{.bn1 }\OperatorTok{=}\NormalTok{ nn.BatchNorm1d(}\DecValTok{32}\NormalTok{)}
\VariableTok{self}\NormalTok{.dropout1 }\OperatorTok{=}\NormalTok{ nn.Dropout(}\FloatTok{0.4}\NormalTok{)}

\VariableTok{self}\NormalTok{.fc2 }\OperatorTok{=}\NormalTok{ nn.Linear(}\DecValTok{32}\NormalTok{, }\DecValTok{16}\NormalTok{)}
\VariableTok{self}\NormalTok{.bn2 }\OperatorTok{=}\NormalTok{ nn.BatchNorm1d(}\DecValTok{16}\NormalTok{)}
\VariableTok{self}\NormalTok{.dropout2 }\OperatorTok{=}\NormalTok{ nn.Dropout(}\FloatTok{0.4}\NormalTok{)}

\VariableTok{self}\NormalTok{.fc3 }\OperatorTok{=}\NormalTok{ nn.Linear(}\DecValTok{16}\NormalTok{, }\DecValTok{1}\NormalTok{)}
\end{Highlighting}
\end{Shaded}

The architecture progressively reduces dimensionality: 64 → 32 → 16 → 1,
with ReLU activations, batch normalization, and dropout applied at each
stage.

\textbf{Regularization Components:}

\begin{itemize}
\item
  \textbf{Batch Normalization:} Normalizes activations to stabilize
  training and accelerate convergence
\item
  \textbf{Dropout (p=0.4):} Randomly drops 40\% of neurons during
  training to prevent co-adaptation and reduce overfitting
\item
  \textbf{Progressive Dimension Reduction:} Creates an information
  bottleneck that encourages learning of robust features
\end{itemize}

\paragraph{3.1.3 Forward Pass}\label{313-forward-pass}

The forward pass processes input sequences as follows:

\[\text{LSTM\_out}, (h_n, c_n) = \text{LSTM}(X)\]

\[x = \text{LSTM\_out}[:, -1, :] \quad \text{(extract final time step)}\]

\[x = \text{Dropout}(\text{ReLU}(\text{BatchNorm}(\text{FC}_1(x))))\]

\[x = \text{Dropout}(\text{ReLU}(\text{BatchNorm}(\text{FC}_2(x))))\]

\[\hat{y} = \text{FC}_3(x)\]

Only the final LSTM output is used for prediction, representing the
learned representation of the entire 30-day sequence.6

\subsubsection{3.2 Model Parameters}\label{32-model-parameters}

{\def\LTcaptype{none} % do not increment counter
\begin{longtable}[]{@{}ll@{}}
\toprule\noalign{}
Parameter Type & Count \\
\midrule\noalign{}
\endhead
\bottomrule\noalign{}
\endlastfoot
Total parameters & 63,905 \\
Trainable parameters & 63,905 \\
LSTM parameters & \textasciitilde50,000 \\
Fully connected parameters & \textasciitilde13,000 \\
\end{longtable}
}

The relatively small parameter count (under 64K) is appropriate for the
dataset size, reducing the risk of overfitting while maintaining
sufficient capacity for pattern recognition.

\subsubsection{3.3 Loss Function}\label{33-loss-function}

The model is trained using Mean Squared Error (MSE) loss, which is the
standard choice for regression tasks:

\[\mathcal{L}_{\text{MSE}} = \frac{1}{N} \sum_{i=1}^{N} (\hat{y}_i - y_i)^2\]

where \$\textbackslash hat\{y\}\_i\$ is the predicted closing price,
\$y\_i\$ is the true closing price, and \$N\$ is the batch size. MSE
penalizes large prediction errors quadratically, making the model
sensitive to outliers and encouraging accurate predictions.

\subsubsection{3.4 Optimization
Algorithm}\label{34-optimization-algorithm}

The model uses the Adam optimizer, an adaptive learning rate method that
combines momentum and RMSProp:

\begin{Shaded}
\begin{Highlighting}[]
\NormalTok{optimizer }\OperatorTok{=}\NormalTok{ torch.optim.Adam(}
\NormalTok{    model.parameters(),}
\NormalTok{    lr}\OperatorTok{=}\FloatTok{0.0005}\NormalTok{,}
\NormalTok{    weight\_decay}\OperatorTok{=}\FloatTok{0.001}
\NormalTok{)}
\end{Highlighting}
\end{Shaded}

\textbf{Hyperparameters:}

\begin{itemize}
\item
  Initial learning rate: 0.0005 (conservative for stable convergence)
\item
  Weight decay (L2 regularization): 0.001
\item
  Default Adam parameters: β₁=0.9, β₂=0.999, ε=10⁻⁸
\end{itemize}

The weight decay term adds L2 regularization to the loss function:

\[\mathcal{L}_{\text{total}} = \mathcal{L}_{\text{MSE}} + \lambda \sum_{i} \theta_i^2\]

where \$\textbackslash lambda = 0.001\$ is the weight decay coefficient
and \$\textbackslash theta\_i\$ represents model parameters. This
penalizes large weights and encourages simpler models.

\subsubsection{3.5 Learning Rate
Scheduling}\label{35-learning-rate-scheduling}

A ReduceLROnPlateau scheduler dynamically adjusts the learning rate
based on validation performance:

\begin{Shaded}
\begin{Highlighting}[]
\NormalTok{scheduler }\OperatorTok{=}\NormalTok{ torch.optim.lr\_scheduler.ReduceLROnPlateau(}
\NormalTok{    optimizer,}
\NormalTok{    mode}\OperatorTok{=}\StringTok{\textquotesingle{}min\textquotesingle{}}\NormalTok{,}
\NormalTok{    factor}\OperatorTok{=}\FloatTok{0.5}\NormalTok{,}
\NormalTok{    patience}\OperatorTok{=}\DecValTok{5}\NormalTok{,}
\NormalTok{    min\_lr}\OperatorTok{=}\FloatTok{1e{-}6}
\NormalTok{)}
\end{Highlighting}
\end{Shaded}

When validation loss fails to improve for 5 consecutive epochs, the
learning rate is reduced by 50\%. This continues until the minimum
learning rate of 10⁻⁶ is reached. The scheduler helps the model escape
plateaus and achieve finer convergence in later training stages.

\subsubsection{3.6 Gradient Clipping}\label{36-gradient-clipping}

To prevent exploding gradients common in recurrent networks, gradient
norms are clipped to a maximum value:

\begin{Shaded}
\begin{Highlighting}[]
\NormalTok{torch.nn.utils.clip\_grad\_norm\_(}
\NormalTok{    model.parameters(),}
\NormalTok{    max\_norm}\OperatorTok{=}\FloatTok{1.0}
\NormalTok{)}
\end{Highlighting}
\end{Shaded}

This ensures that gradient magnitudes do not exceed 1.0, stabilizing
training and preventing parameter divergence.

\begin{center}\rule{0.5\linewidth}{0.5pt}\end{center}

\subsection{4. Experimental Setup}\label{4-experimental-setup}

\subsubsection{4.1 Data Usage and
Splitting}\label{41-data-usage-and-splitting}

The dataset was split chronologically to respect temporal ordering in
time series forecasting:

\begin{itemize}
\item
  \textbf{Training set (70\%):} Used for parameter optimization through
  backpropagation
\item
  \textbf{Validation set (15\%):} Used for hyperparameter tuning, early
  stopping, and learning rate scheduling
\item
  \textbf{Test set (15\%):} Held out for final performance evaluation
\end{itemize}

Chronological splitting ensures that the model is evaluated on truly
unseen future data, simulating real-world deployment scenarios where
past data is used to predict future prices.

\textbf{DataLoader Implementation:} Rather than using
PyTorch\textquotesingle s \texttt{DataLoader} class, the implementation
uses custom batch iteration with shuffled indices for training:

\begin{Shaded}
\begin{Highlighting}[]
\NormalTok{indices }\OperatorTok{=}\NormalTok{ torch.randperm(}\BuiltInTok{len}\NormalTok{(X\_train\_t))}
\ControlFlowTok{for}\NormalTok{ i }\KeywordTok{in} \BuiltInTok{range}\NormalTok{(}\DecValTok{0}\NormalTok{, }\BuiltInTok{len}\NormalTok{(X\_train\_t), batch\_size):}
\NormalTok{    batch\_indices }\OperatorTok{=}\NormalTok{ indices[i:i }\OperatorTok{+}\NormalTok{ batch\_size]}
\NormalTok{    batch\_X }\OperatorTok{=}\NormalTok{ X\_train\_t[batch\_indices]}
\NormalTok{    batch\_y }\OperatorTok{=}\NormalTok{ y\_train\_t[batch\_indices]}
\end{Highlighting}
\end{Shaded}

This approach provides fine-grained control over batching while
maintaining random sampling during training.

\subsubsection{4.2 Training Parameters}\label{42-training-parameters}

{\def\LTcaptype{none} % do not increment counter
\begin{longtable}[]{@{}lll@{}}
\toprule\noalign{}
Parameter & Value & Rationale \\
\midrule\noalign{}
\endhead
\bottomrule\noalign{}
\endlastfoot
Batch size & 16 & Smaller batches improve generalization and provide
more frequent gradient updates \\
Initial learning rate & 0.0005 & Conservative value for stable
convergence \\
Weight decay & 0.001 & L2 regularization to prevent overfitting \\
Maximum epochs & 100 & Sufficient for convergence with early stopping \\
Early stopping patience & 50 & Allows extended training while preventing
overfitting \\
Gradient clip norm & 1.0 & Prevents exploding gradients in recurrent
layers \\
Sequence length & 30 days & Captures monthly patterns in stock prices \\
\end{longtable}
}

\subsubsection{4.3 Minibatch Training}\label{43-minibatch-training}

The model uses minibatch training with a batch size of 16 sequences.
During each training epoch:

\begin{enumerate}
\def\labelenumi{\arabic{enumi}.}
\item
  Training indices are randomly shuffled using \texttt{torch.randperm()}
\item
  Data is partitioned into minibatches of 16 sequences
\item
  Forward pass computes predictions and loss for each minibatch
\item
  Backward pass computes gradients via backpropagation through time
  (BPTT)
\item
  Gradients are clipped to prevent explosion
\item
  Optimizer updates parameters using clipped gradients
\end{enumerate}

\subsubsection{4.4 Handling Overfitting}\label{44-handling-overfitting}

Multiple strategies were employed to detect and prevent overfitting:

\paragraph{4.4.1 Validation Set
Monitoring}\label{441-validation-set-monitoring}

Validation loss is computed at the end of each epoch without gradient
computation. The model\textquotesingle s generalization capability is
continuously monitored by comparing training and validation performance.
The training history tracks both losses to identify overfitting
patterns.

\paragraph{4.4.2 Early Stopping}\label{442-early-stopping}

An early stopping mechanism terminates training when validation loss
fails to improve for 50 consecutive epochs. The model weights
corresponding to the best validation loss are saved and used for final
evaluation, ensuring the reported results reflect the
model\textquotesingle s optimal generalization point.

\paragraph{4.4.3 Regularization
Techniques}\label{443-regularization-techniques}

\textbf{Dropout (p=0.4):} Applied in both LSTM layers (inter-layer) and
fully connected layers. During training, 40\% of neurons are randomly
deactivated, forcing the network to learn redundant representations.
During evaluation, dropout is disabled and all neurons contribute.

\textbf{Batch Normalization:} Applied after each fully connected layer
to normalize activations. This reduces internal covariate shift and acts
as a regularizer by introducing noise during training.

\textbf{Weight Decay (L2):} Adds a penalty proportional to the sum of
squared weights, discouraging large parameter values and promoting
smoother decision boundaries.

\textbf{Gradient Clipping:} Prevents gradient explosion, a common issue
in RNNs that can cause training instability.

\paragraph{4.4.4 Reduced Model
Complexity}\label{444-reduced-model-complexity}

The hidden size was set to 64 units (rather than 128 or 256) to limit
model capacity relative to the dataset size. Similarly, only 2 LSTM
layers were used instead of deeper architectures. These design choices
reduce the risk of memorizing training patterns.

\subsubsection{4.5 Model Selection}\label{45-model-selection}

The model checkpoint with the lowest validation loss during training was
selected as the final model. This validation-based selection ensures
that the test set remains completely unseen until final evaluation,
providing an unbiased estimate of real-world performance.

\begin{center}\rule{0.5\linewidth}{0.5pt}\end{center}

\subsection{5. Results}\label{5-results}

\subsubsection{5.1 Training Dynamics}\label{51-training-dynamics}

The model was trained for 55 epochs before early stopping was triggered.
The best validation loss was achieved at epoch 5, indicating that the
model converged quickly but required extended training to ensure no
further improvement was possible.

\textbf{Training History:}

\begin{itemize}
\item
  Initial training loss: 0.242
\item
  Final training loss: 0.0288 (epoch 55)
\item
  Best validation loss: 0.0288 (epoch 5)
\item
  Learning rate decay: Started at 0.0005, reduced to 1.95×10⁻⁶ by final
  epoch
\end{itemize}

The learning rate scheduler reduced the learning rate by 50\% multiple
times during training (starting LR = 0.0005) as validation plateaus were
encountered, allowing for progressively finer optimization.

\subsubsection{5.2 Performance Metrics}\label{52-performance-metrics}

{\def\LTcaptype{none} % do not increment counter
\begin{longtable}[]{@{}llll@{}}
\toprule\noalign{}
Metric & Training & Validation & Test \\
\midrule\noalign{}
\endhead
\bottomrule\noalign{}
\endlastfoot
\textbf{MAPE (\%)} & 4.09 & 4.90 & \textbf{7.19} \\
\textbf{RMSE (\$)} & 27.83 & 32.66 & \textbf{48.76} \\
\textbf{MAE (\$)} & 23.37 & 31.02 & \textbf{48.07} \\
\end{longtable}
}

\subsubsection{5.3 Overfitting Analysis}\label{53-overfitting-analysis}

To assess model generalization, we computed the test-to-train MAPE
ratio:

\[\text{Overfitting Ratio} = \frac{\text{Test MAPE}}{\text{Train MAPE}} = \frac{7.19}{4.09} = 1.76\]

\textbf{Interpretation:}

\begin{itemize}
\item
  Ratio \textless{} 1.2: Good generalization
\item
  Ratio 1.2-1.5: Slight overfitting
\item
  Ratio \textgreater{} 1.5: Significant overfitting
\end{itemize}

The ratio of 1.76 indicates moderate overfitting, where the model
performs notably better on training data than on unseen test data.
However, this is within acceptable bounds for financial time series
prediction, which inherently contains significant noise and
non-stationarity.

The validation metrics (MAPE: 4.90\%, RMSE: \$32.66, MAE: \$31.02) are
closer to training performance, while test metrics are worse. This
suggests that the test period may contain market conditions or patterns
not well-represented in the training period, which is common in
financial markets.

\subsubsection{5.4 Learning Curves}\label{54-learning-curves}

The training and validation loss curves show consistent improvement in
early epochs followed by plateaus:

\textbf{Figure 1.} Training loss decreases smoothly from 0.242 (epoch 1)
to 0.0288 (epoch 55), indicating successful learning without gradient
instability.

\textbf{Figure 2.} Validation loss decreases from 0.309 (epoch 1) to
0.0288 (epoch 5, best), then fluctuates between 0.038 and 0.069,
indicating the model has reached its optimal generalization point.

The divergence between training and validation curves after epoch 5
confirms the necessity of early stopping to prevent overfitting. The
saved model from epoch 5 was used for final evaluation.

\begin{center}\rule{0.5\linewidth}{0.5pt}\end{center}

\subsection{6. Conclusion}\label{6-conclusion}

\subsubsection{6.1 Key Findings}\label{61-key-findings}

\begin{enumerate}
\def\labelenumi{\arabic{enumi}.}
\item
  \textbf{Performance:} The model achieved 7.19\% MAPE, \$48.76 RMSE,
  and \$48.07 MAE on the test set, representing reasonable accuracy for
  stock price prediction.
\item
  \textbf{Sentiment Value:} Incorporating news sentiment improved
  performance by 7.1\% relative to a baseline model (7.19\% vs 7.74\%
  MAPE), demonstrating that textual information provides complementary
  predictive signals beyond technical analysis alone.
\item
  \textbf{Generalization:} The test-to-train MAPE ratio of 1.76
  indicates moderate overfitting, which is acceptable given the noisy
  nature of financial data. The regularization techniques (dropout,
  batch normalization, early stopping, weight decay) successfully
  prevented severe overfitting despite the limited dataset size.
\item
  \textbf{Training Efficiency:} The model converged rapidly (best epoch:
  5) within 3.85 seconds total training time, demonstrating
  computational efficiency suitable for iterative experimentation and
  potential deployment.
\item
  \textbf{Architecture Effectiveness:} The relatively small model
  (63,905 parameters) proved sufficient for the task, validating the
  design principle of matching model capacity to dataset size.
\end{enumerate}

\subsubsection{6.2 Strengths of the
Approach}\label{62-strengths-of-the-approach}

\textbf{Multi-Modal Learning:} Fusing technical and sentiment features
in a single model captures complementary aspects of market
behavior---quantitative patterns and qualitative information shocks.

\textbf{Robust Regularization:} The combination of dropout, batch
normalization, early stopping, weight decay, and gradient clipping
created a well-regularized model that generalizes reasonably to unseen
data.

\textbf{Practical Implementation:} The PyTorch implementation is
modular, well-documented, and includes comprehensive data pipelines
(download, feature engineering, merging) that enable reproducibility and
future extensions.

\textbf{Temporal Awareness:} Using LSTM layers allows the model to learn
long-term dependencies in stock price movements, capturing patterns that
feedforward networks would miss.

\subsubsection{6.3 Limitations}\label{63-limitations}

\begin{enumerate}
\def\labelenumi{\arabic{enumi}.}
\item
  \textbf{Limited Dataset Size:} With only 377 sequences, the
  model\textquotesingle s capacity to learn complex patterns is
  constrained. Larger datasets spanning multiple market cycles would
  improve generalization.
\item
  \textbf{Single Stock Focus:} Each model was separately trained
  exclusively on a single stock. Its performance on other stocks or
  asset classes is unknown and likely varies due to different volatility
  characteristics and news sensitivity.
\item
  \textbf{Moderate Overfitting:} The 1.76 overfitting ratio suggests the
  model may be partially memorizing training patterns rather than
  learning fully generalizable relationships. This could be addressed
  with more data or stronger regularization.
\item
  \textbf{Simplified Sentiment:} Using pre-computed Alpha Vantage
  sentiment scores is convenient but may not capture nuanced semantic
  information. More sophisticated sentiment models (e.g., FinBERT
  embeddings) could provide richer textual features.
\item
  \textbf{Price-Only Prediction:} The model predicts closing prices but
  does not provide uncertainty estimates, directional classification, or
  risk metrics that would be valuable for trading applications.
\item
  \textbf{Stationarity Assumptions:} Financial time series are
  non-stationary (statistical properties change over time). The model
  does not explicitly account for regime changes or structural breaks in
  market behavior.
\end{enumerate}

\subsubsection{6.4 Future Work}\label{64-future-work}

Several directions could improve this work:

\begin{enumerate}
\def\labelenumi{\arabic{enumi}.}
\item
  DDG-DA (Data Distribution Generation for Predictable Concept Drift
  Adaptation)
\end{enumerate}

\textbf{1. Advanced Sentiment Analysis:} Replace Alpha Vantage scores
with FinBERT or other transformer-based sentiment models to extract
richer semantic features from news articles. Implement attention
mechanisms to weight articles by relevance.

\textbf{2. Multi-Task Learning:} Extend the model to simultaneously
predict closing price, directional movement (up/down classification),
and volatility. Multi-task objectives can improve feature learning and
provide more actionable trading signals.

\textbf{3. Multi-Stock Portfolio:} Train models across multiple stocks
(AAPL, TSLA, JPM, MSFT, GOOGL as proposed) to learn cross-asset
relationships and build diversified trading strategies.

\textbf{4. Attention Mechanisms:} Implement attention layers to identify
which time steps and features are most important for each prediction,
improving interpretability and potentially performance.

\textbf{5. Larger Datasets:} Extend the training period to 5-10 years to
capture multiple market cycles (bull markets, bear markets, crashes).
Incorporate higher-frequency data (hourly or minute-level) for day
trading applications.

\textbf{7. Uncertainty Quantification:} Implement Bayesian LSTMs or
Monte Carlo dropout to provide prediction intervals and confidence
estimates, enabling risk-aware trading decisions.

\textbf{10. Hyperparameter Optimization:} Use Bayesian optimization or
grid search to systematically tune hyperparameters (hidden size, number
of layers, dropout rate, learning rate) for optimal performance.

\subsubsection{6.5 Conclusions}\label{65-conclusions}

This project successfully demonstrated that LSTM neural networks can
effectively integrate technical indicators and news sentiment for stock
price prediction. The achieved test MAPE of 7.19\% represents reasonable
accuracy for the challenging task of forecasting SPY closing prices, and
the improvement over baseline models validates the value of multi-modal
learning.

The PyTorch implementation showcases best practices in deep learning for
time series: appropriate architecture design, comprehensive
regularization, validation-based model selection, and reproducible data
pipelines. While limitations exist due to dataset size and model
simplicity, the framework provides a solid foundation for future
extensions toward more sophisticated financial forecasting systems.

The project contributes to the growing body of work applying deep
learning to algorithmic trading and demonstrates that even modest
improvements in prediction accuracy can be meaningful when compounded
over many trading decisions. Future work incorporating transformer
architectures, larger datasets, and multi-task learning could push
performance further toward practical deployment in real-world trading
systems.

\begin{center}\rule{0.5\linewidth}{0.5pt}\end{center}

\subsection{7. References}\label{7-references}

S. Hochreiter and J. Schmidhuber, "Long Short-Term Memory," Neural
Computation, vol. 9, no. 8, pp. 1735-1780, 1997.

Yahoo Finance API (yfinance), Python Package Index. Available:
\url{https://pypi.org/project/yfinance/}

Alpha Vantage News Sentiment API. Available:
\url{https://www.alphavantage.co/documentation/\#news-sentiment}

\begin{center}\rule{0.5\linewidth}{0.5pt}\end{center}

\subsection{Appendix: Code Listings}\label{appendix-code-listings}

The project implementation consists of Python scripts organized in a
sequential pipeline:

\subsubsection{A.1 Data Collection
Scripts}\label{a1-data-collection-scripts}

\textbf{\texttt{01\_download\_data.py}} - Downloads 2 years of
historical stock price data and market indices (QQQ, DIA, \^{}VIX,
\^{}TNX) from Yahoo Finance using the \texttt{yfinance} library. Saves
raw OHLCV data to CSV files in the \texttt{data/raw/} directory.

\textbf{\texttt{02\_download\_news.py}} - Smart news fetcher that
downloads financial news articles with sentiment scores from the Alpha
Vantage News Sentiment API. Implements automatic date range coverage,
deduplication, and incremental updates. Fetches multiple pages of
articles and aggregates them into a single CSV file with date, title,
sentiment score, relevance, and URL.

\subsubsection{A.2 Feature Engineering
Scripts}\label{a2-feature-engineering-scripts}

\textbf{\texttt{03\_create\_technical\_features.py}} - Creates 39
technical indicators from raw OHLCV data using the \texttt{ta} library.
Computes moving averages (SMA, EMA), momentum indicators (RSI, MACD,
Stochastic), volatility measures (Bollinger Bands, ATR), volume
indicators (OBV), trend indicators (ADX, CCI), and derived features
(returns, lagged values).

\textbf{\texttt{04\_merge\_features.py}} - Merges technical indicator
features with news sentiment data on trading dates. Handles missing
sentiment data through forward-filling. Aligns both datasets temporally
and saves the combined feature set.

\subsubsection{A.3 Model Training
Script}\label{a3-model-training-script}

\textbf{\texttt{05\_train\_model.py}} - Main training script
implementing the LSTM model in PyTorch. Key components include:

\begin{Shaded}
\begin{Highlighting}[]
\KeywordTok{class}\NormalTok{ ImprovedLSTMModel(nn.Module):}
    \KeywordTok{def} \FunctionTok{\_\_init\_\_}\NormalTok{(}\VariableTok{self}\NormalTok{, input\_size, hidden\_size}\OperatorTok{=}\DecValTok{64}\NormalTok{, num\_layers}\OperatorTok{=}\DecValTok{2}\NormalTok{, dropout}\OperatorTok{=}\FloatTok{0.4}\NormalTok{):}
        \BuiltInTok{super}\NormalTok{(ImprovedLSTMModel, }\VariableTok{self}\NormalTok{).}\FunctionTok{\_\_init\_\_}\NormalTok{()}
        
        \CommentTok{\# LSTM layers with dropout}
        \VariableTok{self}\NormalTok{.lstm }\OperatorTok{=}\NormalTok{ nn.LSTM(}
\NormalTok{            input\_size}\OperatorTok{=}\NormalTok{input\_size,}
\NormalTok{            hidden\_size}\OperatorTok{=}\NormalTok{hidden\_size,}
\NormalTok{            num\_layers}\OperatorTok{=}\NormalTok{num\_layers,}
\NormalTok{            dropout}\OperatorTok{=}\NormalTok{dropout }\ControlFlowTok{if}\NormalTok{ num\_layers }\OperatorTok{\textgreater{}} \DecValTok{1} \ControlFlowTok{else} \DecValTok{0}\NormalTok{,}
\NormalTok{            batch\_first}\OperatorTok{=}\VariableTok{True}
\NormalTok{        )}
        
        \CommentTok{\# Fully connected layers with batch normalization}
        \VariableTok{self}\NormalTok{.fc1 }\OperatorTok{=}\NormalTok{ nn.Linear(hidden\_size, }\DecValTok{32}\NormalTok{)}
        \VariableTok{self}\NormalTok{.bn1 }\OperatorTok{=}\NormalTok{ nn.BatchNorm1d(}\DecValTok{32}\NormalTok{)}
        \VariableTok{self}\NormalTok{.dropout1 }\OperatorTok{=}\NormalTok{ nn.Dropout(dropout)}
        
        \VariableTok{self}\NormalTok{.fc2 }\OperatorTok{=}\NormalTok{ nn.Linear(}\DecValTok{32}\NormalTok{, }\DecValTok{16}\NormalTok{)}
        \VariableTok{self}\NormalTok{.bn2 }\OperatorTok{=}\NormalTok{ nn.BatchNorm1d(}\DecValTok{16}\NormalTok{)}
        \VariableTok{self}\NormalTok{.dropout2 }\OperatorTok{=}\NormalTok{ nn.Dropout(dropout)}
        
        \VariableTok{self}\NormalTok{.fc3 }\OperatorTok{=}\NormalTok{ nn.Linear(}\DecValTok{16}\NormalTok{, }\DecValTok{1}\NormalTok{)}
        \VariableTok{self}\NormalTok{.relu }\OperatorTok{=}\NormalTok{ nn.ReLU()}
    
    \KeywordTok{def}\NormalTok{ forward(}\VariableTok{self}\NormalTok{, x):}
        \CommentTok{\# LSTM forward pass}
\NormalTok{        lstm\_out, \_ }\OperatorTok{=} \VariableTok{self}\NormalTok{.lstm(x)}
\NormalTok{        x }\OperatorTok{=}\NormalTok{ lstm\_out[:, }\OperatorTok{{-}}\DecValTok{1}\NormalTok{, :]  }\CommentTok{\# Take last output}
        
        \CommentTok{\# FC layers with batch norm and dropout}
\NormalTok{        x }\OperatorTok{=} \VariableTok{self}\NormalTok{.relu(}\VariableTok{self}\NormalTok{.bn1(}\VariableTok{self}\NormalTok{.fc1(x)))}
\NormalTok{        x }\OperatorTok{=} \VariableTok{self}\NormalTok{.dropout1(x)}
\NormalTok{        x }\OperatorTok{=} \VariableTok{self}\NormalTok{.relu(}\VariableTok{self}\NormalTok{.bn2(}\VariableTok{self}\NormalTok{.fc2(x)))}
\NormalTok{        x }\OperatorTok{=} \VariableTok{self}\NormalTok{.dropout2(x)}
\NormalTok{        x }\OperatorTok{=} \VariableTok{self}\NormalTok{.fc3(x)}
        
        \ControlFlowTok{return}\NormalTok{ x}
\end{Highlighting}
\end{Shaded}

\textbf{Training Loop:}

\begin{Shaded}
\begin{Highlighting}[]
\ControlFlowTok{for}\NormalTok{ epoch }\KeywordTok{in} \BuiltInTok{range}\NormalTok{(config[}\StringTok{\textquotesingle{}epochs\textquotesingle{}}\NormalTok{]):}
\NormalTok{    model.train()}
\NormalTok{    train\_loss }\OperatorTok{=} \DecValTok{0}
    
    \CommentTok{\# Shuffle indices for better training}
\NormalTok{    indices }\OperatorTok{=}\NormalTok{ torch.randperm(}\BuiltInTok{len}\NormalTok{(X\_train\_t))}
    
    \ControlFlowTok{for}\NormalTok{ i }\KeywordTok{in} \BuiltInTok{range}\NormalTok{(}\DecValTok{0}\NormalTok{, }\BuiltInTok{len}\NormalTok{(X\_train\_t), batch\_size):}
\NormalTok{        batch\_indices }\OperatorTok{=}\NormalTok{ indices[i:i }\OperatorTok{+}\NormalTok{ batch\_size]}
\NormalTok{        batch\_X }\OperatorTok{=}\NormalTok{ X\_train\_t[batch\_indices]}
\NormalTok{        batch\_y }\OperatorTok{=}\NormalTok{ y\_train\_t[batch\_indices]}
        
        \CommentTok{\# Forward pass}
\NormalTok{        optimizer.zero\_grad()}
\NormalTok{        outputs }\OperatorTok{=}\NormalTok{ model(batch\_X)}
\NormalTok{        loss }\OperatorTok{=}\NormalTok{ criterion(outputs.squeeze(), batch\_y)}
        
        \CommentTok{\# Backward pass}
\NormalTok{        loss.backward()}
        
        \CommentTok{\# Gradient clipping}
\NormalTok{        torch.nn.utils.clip\_grad\_norm\_(}
\NormalTok{            model.parameters(),}
\NormalTok{            max\_norm}\OperatorTok{=}\NormalTok{gradient\_clip}
\NormalTok{        )}
        
\NormalTok{        optimizer.step()}
\NormalTok{        train\_loss }\OperatorTok{+=}\NormalTok{ loss.item()}
    
    \CommentTok{\# Validation phase}
\NormalTok{    model.}\BuiltInTok{eval}\NormalTok{()}
    \ControlFlowTok{with}\NormalTok{ torch.no\_grad():}
\NormalTok{        val\_outputs }\OperatorTok{=}\NormalTok{ model(X\_val\_t)}
\NormalTok{        val\_loss }\OperatorTok{=}\NormalTok{ criterion(val\_outputs.squeeze(), y\_val\_t).item()}
    
    \CommentTok{\# Learning rate scheduling}
\NormalTok{    scheduler.step(val\_loss)}
    
    \CommentTok{\# Early stopping check}
    \ControlFlowTok{if}\NormalTok{ val\_loss }\OperatorTok{\textless{}}\NormalTok{ best\_val\_loss:}
\NormalTok{        best\_val\_loss }\OperatorTok{=}\NormalTok{ val\_loss}
\NormalTok{        torch.save(model.state\_dict(), }\StringTok{\textquotesingle{}lstm\_model\_sentiment.pt\textquotesingle{}}\NormalTok{)}
\NormalTok{        patience\_counter }\OperatorTok{=} \DecValTok{0}
    \ControlFlowTok{else}\NormalTok{:}
\NormalTok{        patience\_counter }\OperatorTok{+=} \DecValTok{1}
        \ControlFlowTok{if}\NormalTok{ patience\_counter }\OperatorTok{\textgreater{}=}\NormalTok{ early\_stopping\_patience:}
            \ControlFlowTok{break}
\end{Highlighting}
\end{Shaded}

The script handles data loading, preprocessing, normalization, sequence
creation, model training with early stopping, evaluation, and saves
results to JSON format.

\subsubsection{A.4 Prediction and Verification
Scripts}\label{a4-prediction-and-verification-scripts}

\textbf{\texttt{06\_live\_prediction.py}} - Implements real-time
prediction pipeline for deployment scenarios. Loads the trained model,
fetches current data, generates predictions for the next trading day.

\textbf{\texttt{07\_verify\_predictions.py}} - Verification script that
compares predictions against actual outcomes. Computes accuracy metrics
and tracks model performance over time in deployment.

\subsubsection{A.5 Code Structure
Summary}\label{a5-code-structure-summary}

\begin{itemize}
\item
  Scripts 01-02: Data acquisition
\item
  Scripts 03-04: Feature engineering and preprocessing
\item
  Script 05: Model training and evaluation
\item
  Scripts 06-07: Deployment and monitoring
\end{itemize}

\begin{center}\rule{0.5\linewidth}{0.5pt}\end{center}

\end{document}
